\chapter{Introduction}
Introduction


\section{Background and Motivation}
Computer science faculty laboratories are energy-intensive academic spaces because they

commonly contain desktop computers, networking switches and routers, projectors, over-

head lighting circuits, cooling fans, and continuously occupied workstations. Unlike ordi-

nary lecture rooms, laboratories are often used according to irregular schedules that include

taught sessions, open laboratory periods, examinations, maintenance work, and temporary

reassignment between instructors. In many faculties, these spaces remain illuminated and

ventilated even when no occupants are present, because the responsible person may have

left without switching off appliances, or because the schedule indicates that a session should

be active even though the instructor has moved the class elsewhere. The cumulative effect

of such waste across dozens of laboratories and hundreds of weekly sessions represents a

substantial and avoidable electricity cost for the faculty and the university.

Recent smart-building research identifies occupancy information as a central input for

energy-efficient control of lighting and heating, ventilation, and air-conditioning systems.

Chaudhari et al. conducted a comprehensive survey of occupancy detection technologies

for smart buildings using IoT sensors, demonstrating that occupancy-aware automation can

meaningfully reduce energy consumption by ensuring that lighting and HVAC systems op-

erate only when and where occupants are actually present [?]. Song et al. further examined

the impact of occupancy behaviour on building energy efficiency, showing that occupant

behaviour is a primary variable in academic building energy demand, and that detection

and monitoring technologies are essential for closing the gap between designed and actual

building performance [?]. Waghchoure et al. demonstrated a practical classroom automa-

tion system that uses PIR sensing and camera activation to trigger appliance control, con-

firming that low-cost IoT components can be deployed in educational settings to reduce

unnecessary electricity use [?].

GreenLab is proposed as an intelligent energy management system for computer sci-

ence faculty laboratories. The system combines embedded sensing, automated appliance

control, computer-vision occupancy estimation, backend decision logic, role-based moni-

toring, and manual override through a web dashboard. The hardware layer is centred on an

ESP32 microcontroller connected to a light-dependent resistor for natural-light sensing, a

DHT11 temperature and humidity sensor, an HC-SR04 ultrasonic distance sensor for entry

and activity detection, and a 4-channel relay module for controlling lights and fans. The

artificial-intelligence layer uses camera frames and a YOLO-based model to estimate the

number of occupants in the laboratory. The backend layer receives sensor and occupancy

data, applies operational rules, stores audit records, and exposes secure REST endpoints.

The frontend layer presents real-time status, laboratory assignments, manual controls, and

administrative actions through an Angular single-page application dashboard.

The project is motivated by the limitations of purely manual control and purely schedule-

based automation. A schedule may indicate that a laboratory should be active, but it does

not confirm whether occupants are present, whether natural light is sufficient, whether ven-

tilation is required, or whether an instructor needs to override the automatic policy for peda-

gogical reasons. Similarly, passive sensing alone may detect motion without distinguishing

sustained occupancy from transient passage, while camera-based detection may be affected

by lighting conditions, occlusion, privacy restrictions, or processing availability. The re-

viewed literature therefore supports multimodal approaches that combine environmental

sensors, occupancy estimation, and supervisory control for practical smart-building systems

[?, ?]. GreenLab follows this direction by using embedded sensor readings for immediate

environmental awareness and computer vision for improved occupancy interpretation, all

governed by a backend rule engine and role-based dashboard.


\section{Problem Statement}
The problem addressed by GreenLab is the absence of an integrated, faculty-oriented sys-

tem that can reduce unnecessary laboratory energy use while preserving instructor con-

trol, operational accountability, and safe equipment management. In many laboratories

at Mansoura University's Faculty of Engineering, lights and fans are controlled manually

through wall switches, and responsibility for switching appliances off is distributed infor-

mally among instructors, students, and laboratory staff. This arrangement is vulnerable

to forgotten switches, delayed response after class changes, and inconsistent handling of

temporary laboratory transfers. When a laboratory becomes unoccupied or only partially

occupied, appliances may continue operating at the same level because no automated mech-

anism links occupancy, light level, temperature, schedule, and user authority.

The problem is not limited to electricity consumption. It also concerns laboratory us-

ability and governance. Instructors require visibility of the laboratory state during sessions

and need a controlled method to override automatic decisions when teaching conditions

require it, such as dimming lights for a projector demonstration or activating ventilation

before students arrive on a hot day. Administrators require logs, user-account manage-

ment, laboratory scheduling, and approval workflows for assignment changes so that room

transfers do not bypass access control or leave appliances in an unmanaged state.

Hardware faults, communication failures, sensor anomalies, and relay actuation is-

sues must be detectable and auditable so that the system can enter a safe state and alert the

appropriate personnel. Challa et al. demonstrated that combining gas sensing with visual

detection improved HVAC response reliability in classrooms, highlighting the importance

of multimodal sensing and governed control paths [?]. Chaudhari et al. confirmed that

occupancy-aware automation reduces building energy waste, but noted that sensor place-

ment, real-time constraints, and privacy considerations remain practical challenges [?].

The absence of a unified system means that each of these concerns is currently ad-

dressed in isolation, if at all. A standalone motion sensor may switch a light but cannot

enforce role-based override or record audit trails. A web dashboard may display labo-

ratory status but cannot actuate relays without a backend authority. A schedule may de-

fine when a laboratory should be active but cannot confirm that the assigned instructor is

actually present. GreenLab is designed to bridge these gaps by providing an integrated

cyber-physical system that connects embedded sensing, occupancy detection, rule-based

automation, supervised override, and administrative governance in one coherent platform.


\section{Project Aim}
The aim of GreenLab is to design an intelligent energy management system that automati-

cally controls lighting and ventilation in computer science faculty laboratories according to

occupancy, environmental conditions, schedule context, and authorised user input. The sys-

tem is intended to support energy-aware operation, reduce avoidable manual effort, preserve

safety through controlled relay actuation, and provide administrators with clear monitoring

and audit capabilities. This aim is pursued through a four-tier architecture that separates

hardware sensing, AI processing, backend decision-making, and dashboard presentation,

enabling each tier to evolve independently while maintaining well-defined interface con-

tracts.


\section{Project Objectives}
The project objectives define the specific technical goals that GreenLab must achieve. Each

objective is expressed as a full statement of intent with sufficient technical detail to guide

the subsequent requirements, design, and implementation chapters.

\begin{itemize}
  \item Objective 1: Define a layered four-tier architecture with clear interface con-
\end{itemize}
tracts. Establish a layered architecture that separates hardware sensing and actuation,

AI processing, backend decision-making and persistence, and frontend presentation

into four distinct tiers. Each tier communicates through well-defined interfaces so

that changes in one tier, such as replacing the YOLO model or redesigning the dash-

board, do not require modifications to the other tiers. The hardware tier uses the

ESP32 microcontroller to read sensors and control relays, communicating with the

backend through REST endpoints. The AI tier runs a Python YOLO service that re-

ceives camera frames and returns occupancy counts. The backend tier implements

an ASP.NET Core REST API with authentication, business rules, and SQL Server

persistence. The frontend tier provides an Angular single-page application for mon-

itoring and control. This separation is essential for maintainability, testability, and

independent evolution of each system component.

\begin{itemize}
  \item Objective 2: Specify functional and non-functional requirements in formal IEEE 830
\end{itemize}
style. Specify functional and non-functional requirements in a formal IEEE 830 style

across nine functional areas: system initialisation, entry and exit detection, occu-

pancy detection, lighting control, ventilation control, monitoring, manual override,

fault handling, and authentication and access control. Each functional requirement

must describe one specific, observable, verifiable behaviour of the system, with prior-

ity and source annotations. Non-functional requirements must address performance,

reliability, security, usability, maintainability, and scalability with measurable targets

where possible. This formal specification provides the contractual basis for design

decisions and future verification, ensuring that every feature in the system traces to a

documented requirement.

\begin{itemize}
  \item Objective 3: Design occupancy detection using multimodal input. Design oc-
\end{itemize}
cupancy detection using multimodal input that combines ESP32 sensor readings,

specifically the HC-SR04 ultrasonic distance sensor for motion and entry detection,

the DHT11 temperature sensor for environmental context, and the LDR for light-level

awareness, with YOLO-based visual counting from an IP camera stream. Multimodal

detection is necessary because no single sensor provides reliable occupancy informa-

tion under all conditions: ultrasonic sensors may produce false positives from non-

human movement, cameras may be affected by lighting and occlusion, and environ-

mental sensors alone cannot confirm human presence. The combination of embedded

sensor activity and camera-derived person counts, filtered by confidence thresholds

and temporal consistency rules, provides a more robust occupancy estimate than any

single modality.

\begin{itemize}
  \item Objective 4: Design lighting control with a decision matrix. Design lighting con-
\end{itemize}
trol with a decision matrix that uses LDR level, occupancy state, and override state

as inputs. The decision matrix must cover all relevant combinations of natural-light

zones (dark, dim, bright), occupancy states (occupied, vacant, unknown), and over-

ride conditions (no override, instructor override, administrator override) to produce

the appropriate light relay command. This approach replaces ad-hoc threshold logic

with a structured, auditable rule table that can be reviewed, configured, and extended

without modifying firmware code. The decision matrix is evaluated by the backend

rule engine, ensuring that lighting commands are authorised, logged, and consistent

with the current system state.

\begin{itemize}
  \item Objective 5: Design ventilation control with a fan activation table. Design venti-
\end{itemize}
lation control with a fan activation table that uses temperature zones and occupancy

classification as inputs. The table must define how many fan relay channels are ac-

tivated and at what speed setting, based on the current temperature zone (cool, com-

fortable, warm, hot) and the occupancy classification (empty, low, moderate, high).

This tabular approach ensures that ventilation decisions are explicit, consistent, and

traceable, and that fan behaviour can be tuned per laboratory without rewriting em-

bedded logic. The ventilation control must also respect override commands, fault

states, and safety rules that take priority over automatic fan activation.

\begin{itemize}
  \item Objective 6: Implement a secure ASP.NET Core REST API. Implement a se-
\end{itemize}
cure ASP.NET Core REST API with authentication, role-based access control, input

validation, and audit logging. The API must serve as the single authority for all state-

changing actions, including sensor ingestion, occupancy updates, relay commands,

manual overrides, transfer requests, user management, and fault reporting. Authen-

tication must use JWT bearer tokens with configurable expiry. Role-based access

must distinguish between instructor, administrator, and service roles, with each role

permitted to access only its authorised endpoints. Input validation must reject mal-

formed or out-of-range payloads before they reach business logic. Audit logging

must record every state-changing request with actor, action, entity, timestamp, and

outcome, creating a tamper-resistant operational record for administrative review.

\begin{itemize}
  \item Objective 7: Implement an Angular dashboard with monitoring, override, schedul-
\end{itemize}
ing, and administration. Implement an Angular single-page application dashboard

that provides real-time monitoring of laboratory status, manual override controls for

authorised instructors and administrators, scheduling views for laboratory sessions,

transfer request forms for laboratory assignment changes, and administrative man-

agement screens for user accounts, logs, and fault review. The dashboard must

present occupancy, sensor values, relay states, active overrides, and fault warnings

in a clear, role-appropriate interface. Instructor views must be restricted to assigned

laboratories, while administrator views must provide access to all laboratories within

their authority. Real-time updates must be delivered through polling or signal-based

mechanisms so that dashboard users see current system state without manual refresh.


\section{Project Timeline}
The GreenLab project spans nine academic months, from October to June, organised into

eight sequential phases. Table 1.1 summarises the phase focus and core deliverable for each

month.

\begin{longtable}{|L{1.8cm}|L{2.0cm}|L{5.0cm}|L{5.5cm}|}
\caption{GreenLab Monthly Project Timeline (October--June)} \label{tab:1-1} \\
\hline
\rowcolor{gray!25}
\textbf{Month} & \textbf{Phase} & \textbf{Focus} & \textbf{Core Deliverable} \\
\hline
\endfirsthead
\hline
\rowcolor{gray!25}
\textbf{Month} & \textbf{Phase} & \textbf{Focus} & \textbf{Core Deliverable} \\
\hline
\endhead
\hline
\endfoot
\hline
\endlastfoot
October & Phase 1 & Literature Review \& Requirements Foundation & Literature review document; draft requirements specification \\
\hline
November & Phase 2 & Use Cases \& Feature Definition & Use case diagram; feature definition document \\
\hline
December & Phase 3 & Architecture \& Database Schema & Four-tier architecture diagram; 11-table database schema; API specification \\
\hline
January & Phase 4 & ESP32 Firmware Development & Working firmware with sensor polling and relay control \\
\hline
February & Phase 5 & ASP.NET Core Backend Development & Functional REST API with automation engine and audit middleware \\
\hline
March & Phase 6 & Python YOLO Occupancy Service & Working YOLO service integrated with backend \\
\hline
April & Phase 7 & Angular Dashboard \& Full-System Integration & Working dashboard; end-to-end integration test log \\
\hline
May & Phase 8a & Report Writing \& Consolidation & Draft Chapters 3--6 of the LaTeX report \\
\hline
June & Phase 8b & Final Compilation \& Submission & Compiled seven-chapter PDF and submission package \\
\hline
\end{longtable}

\section{Scope}
The project scope includes the design of the GreenLab system architecture, requirements

specification, operational features, workflow modelling, diagrams, and implementation nar-

rative. It encompasses automatic lighting control based on occupancy and natural-light

conditions, automatic ventilation control based on occupancy and temperature context,

continuous monitoring of laboratory state through a web dashboard, manual override for

authorised instructors and administrators, authenticated dashboard access with role-based

permissions, laboratory assignment change handling with approval workflows, fault detec-

tion and reporting, and audit logging of all state-changing actions. The scope also includes

a literature review focused on IoT energy management, classroom and laboratory automa-

tion, YOLO-based occupancy detection, ESP32-based control systems, and web dashboard

architectures for IoT systems.

The scope explicitly excludes a dedicated testing and evaluation chapter. No mea-

sured energy savings percentages, experimental accuracy values for YOLO detection on

GreenLab laboratory footage, measured latency results, or user-satisfaction study findings

are reported because such measurements have not been supplied as verified project outputs.

Where research papers report their own metrics, those metrics are attributed to the original

sources and are not presented as GreenLab results [?, ?]. The scope also excludes physical

installation of hardware in production laboratories, final deployment of the AI service on

production servers, and any regulatory compliance assessment beyond the institutional IT

policy of Mansoura University.


\section{Project Organisation}
This project is organised into seven chapters. Chapter 1 introduces the project background,

problem statement, aim, objectives, timeline, and scope. Chapter 2 reviews relevant re-

search papers and project references, providing an explicit analysis of each source's rele-

vance and a comparative analysis that identifies the literature gap addressed by GreenLab.

Chapter 3 presents the system analysis, including stakeholder analysis, user requirements,

use case descriptions, functional requirements in IEEE 830 style across nine areas with a

total of 135 requirements, non-functional requirements across six quality dimensions, per-

formance budgets, technical constraints, and a requirements traceability matrix. Chapter 4

describes the twelve system features and their operational workflows, including sequence

diagrams for four key workflows and a comparison of operating modes. Chapter 5 presents

the detailed system design, covering the four-tier architecture, hardware design with sensor

specifications, database design with eleven tables, REST API design with eighteen end-

points, AI processing design, lighting control decision matrix, fan activation levels, fron-

tend design, and security design. Chapter 6 provides a grounded implementation narrative

with representative code snippets for the ESP32 firmware, ASP.NET Core backend, Python

YOLO service, and Angular dashboard. Chapter 7 concludes the report, summarises the

project contributions, acknowledges limitations, and identifies future work directions.


